\documentclass[11pt]{article}
\usepackage[T1]{fontenc}\usepackage{lmodern,amsmath,amssymb,amsthm,booktabs,geometry}
\geometry{margin=.92in}\setlength{\emergencystretch}{2em}
\newtheorem{theorem}{Theorem}\newtheorem{lemma}{Lemma}
\newcommand{\dd}{\,d}\newcommand{\E}{\mathbb E}
\title{Critical strong entry: two comparison profiles and a leading saturation bound}
\author{Standalone analytical increment}\date{Version 1, November 2026}
\begin{document}\maketitle
\begin{abstract}
At the critical nominal clock the zero-noise Q1 profile has an explicit
rational density. At positive noise a nonvanishing Gaussian gate term
changes the limiting density and shifts its lower support edge to the
resident angle. A one-dimensional quadrature gives the corrected
probe-active fraction. In the leading system at zero weak mass,
coordinate-ordered input and output profiles retain a positive output
gap through strong saturation. We distinguish this finite-amplitude
estimate from linear instability at the resident and from the open
transfer to the trained flow.
\end{abstract}

\section{Models, source domains and powers}\label{an04cfsv1:scope}
Fix \(\rho\in[13/20,7/10]\), \(\lambda=\rho^2,\ \omega=(1-\lambda)/2\),
\(S_0=cq^{1/\omega},\ c,A>0,\ \ell=\log(1/q)\), and
\(T_A=\log(A/S_0)\). The nominal mass scale refers to \(S_0e^{T_A}=A\).
Write \(b=(c/A)^\omega,\ K(t)=\phi(t)+t\Phi(t)\);
\(a=\rho K(\rho a)>0,\ P=\Phi(\rho a),\
\alpha_{\rm sh}=\lambda P/2,\ d=1/2-\alpha_{\rm sh},\ s=d/\omega\).
The initial label law is \(d\beta/(2\pi)\) with aligned input/output and
uniform squared radius \(S_0\).
\begin{itemize}
\item TAIL \texttt{outer, zerofield, quadrants} applies to the auxiliary
point-mass target-only data model, each open quadrant, all times and
all \(S_0>0\). Its \texttt{exponent} corollary adds
\((S_0/A)^\omega/q\to0\) at \(T_A\).
\item PROFILE \texttt{angle, coordinate, root} concerns positive-noise
target-only dynamics, small q, ratios on compact subsets of \((0,1)\).
Its fixed-chart expansion uses bounded \(\theta/q\). Its upper-half
Q1 arrival law uses \(\beta\in[\pi/4,\pi/2]\), and the entry theorem
for all Q1 weights and the cutoff requires
\(T\ge(2/\omega)\ell+C_1,\ q^2T\le1\).
\item Q2 crossing/transit and CRA \texttt{outer} refer to compact
interior Q2 ancestry in the comparison, with a strict exponent gap for
CRA's pre-crossing result. An endpoint or an original trajectory
requires separate estimates.
\item WIT \texttt{popsystem} and P \texttt{leading, resident, spectra}
apply within the defined leading system on chart-valid finite
intervals with adequate moments. We use positive compactly supported
strong measures, \(N=0,\ 0<M\le1\). The spectrum is at a coherent
resident and separates centered from coherent perturbations.
LOC's local gain additionally needs a small shape, an entered weak
clock and controlled mean, feedback and field errors.
\item For the initialized trained model P's early comparison and
TAIL weighted ancestry use the early alignment guard. ARR main uses
\(3/5\le\rho\le3/4,\ q\le1/20,\ S_0<\delta\le1/100\) and
\(R\ge1,\ qR\le\lambda/2\). OWPT, SMCP and MOMC are ten-power
statements with an all-right-half chart before saturation. Their
specific clock tests are audited in Section~\ref{an04cfsv1:original}.
\end{itemize}

The critical coefficient satisfies
\[
 800/231\le1/\omega\le200/51<4<2/\lambda,\qquad
 1/\omega<2L_-,\quad L_-=2+169/231 .
\]
Here \(2/\lambda\ge200/49>4\). Thus compact interior Q2 labels
are still left-facing in the comparison at \(T_A\); the gap from the
crossing coefficient is uniform. With the logarithmic convention
\(\kappa_q=\log(1/S_0)/\ell\), the exact sign is
\(\kappa_q=1/\omega-\log c/\ell=1/\omega+\log c/\log q\).
In particular it is the limiting exponent that depends on \(\rho\).
Direct substitution into the \emph{candidate} laws gives
\[
 \epsilon(T_A)=q^{-s}e^{-dT_A}=(c/A)^d,\qquad
 E_{\rm loc}(1/\omega)=\frac{1-P}{1-\lambda}
                     -\frac{1-\lambda P}{1-\lambda}=-P .
\]
Formally \(q^{E_{\rm loc}(\kappa_q)}=c^{(1-P)/2}q^{-P}\).
These powers do not supply a transferred tail or justify continuing a
local multiplier past its radius of validity.

\section{Q1 comparison at the critical clock}\label{an04cfsv1:profiles}
Write \(x=\theta/q\). Densities below are radial masses per unit \(x\);
the denominator for activity is the total Q1-ancestry mass.
Define
\[
 g_b(h)=\frac{A}{2\pi b}\left(1+(h/b)^2\right)^{-2}\quad(h>0),
 \qquad
 F_b(l)=\frac4\pi\int_0^{\max(l,0)/b}\frac{\dd t}{(1+t^2)^2}.
\]
\begin{theorem}[Zero-noise and positive-noise laws]\label{an04cfsv1:laws}
For \(\beta\) on any fixed compact subset of \((0,\pi/2)\),
put \(h=b\tan\beta\). In the zero-noise target-only model,
\(x(T_A)=h+o(1)\); its limiting density is \(g_b(x)\) for \(x>0\),
and its active fraction is \(F_b(\zeta)\).

For positive-noise target-only flow define, for \(w>a\),
\[
 B(w)=\omega w-\tfrac\rho2K(-\rho w)
     =\tfrac12[w-\rho K(\rho w)],\qquad
 H(w)=w\exp\left\{-\omega\int_w^\infty
       \left(\frac1{B(r)}-\frac1{\omega r}\right)\dd r\right\}.
\]
Then
\begin{align}
 x(T_A)&=H^{-1}(h)+o(1),&
 m(T_A,\beta)&= A\cos^2\beta\,(1+o(1)).
 \label{an04cfsv1:positive}
\end{align}
The same mass law holds in the zero-noise model. For positive noise the
limiting density is \(g_b(H(x))H'(x)\) on \(x>a\), zero on \(x\le a\).
The limiting Q1 probe-active fraction is
\begin{equation}\label{an04cfsv1:active}
 F_+(c,\zeta;\rho,A)=
 \begin{cases}0,&\text{if }\zeta\le a,\\ F_b(H(\zeta)),&\text{if }\zeta>a.\end{cases}
\end{equation}
Both Q1 masses converge to \(A/8\). Angle and relative weight errors
are uniform on the specified label compact, the ratio box, and compact
positive ranges of \(c,A\). Fraction convergence is uniform also for
\(\zeta\) on any fixed bounded interval.
\end{theorem}
\begin{proof}
TAIL quadrants gives \(\tan\theta=e^{-\omega t}\tan\beta\) and
\(m=S_0e^t(\cos^2\beta+e^{-2\omega t}\sin^2\beta)\).
At \(T_A\), \(e^{-\omega T_A}=qb\). Changing variables from
\(\beta\) to \(h\) in \(A\cos^2\beta\,d\beta/(2\pi)\) gives \(g_b\).

For the second law set \(w=\tan\theta/q,\ w_0=\tan\beta/q\),
\(D=(1+q^2w^2)^{1/2}\). PROFILE's exact angle equation gives
\(w'=-B_q(w)\), where
\[
 B_q(w)=\omega w+\tfrac12qwD K(-1/(qD))
                     -\tfrac\rho2D K(-\rho w/D).
\]
On bounded \(w\)-intervals \(B_q\to B\) uniformly. Since
\(B'(w)=(1-\lambda\Phi(\rho w))/2\ge\omega>0\),
\(B(a)=0\) is simple, with \(B'(a)=d\). For \(w>a\) the integrand
in \(H\) is strictly positive, integrable at infinity by the Gaussian
tail, and has a divergent integral at \(a\). Thus \(H(a+)=0\),
\(H(w)/w\to1\) at infinity, and
\(H'/H=\omega/B>0\). This proves invertibility and the strict
inequality \(H(w)<w\).

For fixed \(w>a\), let \(\tau_q(w)\) be its first hit from \(w_0\).
On our label compact \(D\le D_*=\sec(\sup\beta)\) between \(w\) and
\(w_0\). For \(r\ge L\) with L fixed sufficiently large,
\(B_q(r)\ge\omega r/2\) and
\[
 \left|\frac1{B_q(r)}-\frac1{\omega r}\right|
 \le C\frac{K(-\rho_-r/D_*)}{r^2}
    +C\frac{qK(-1/(qD_*))}{r},\qquad r\le w_0 .
\]
The second term integrates to \(o(1)\) up to \(w_0=O(q^{-1})\);
the first has uniformly vanishing integral as \(L\to\infty\).
On \([w,L]\) positivity and convergence hold for small q. Separating
variables therefore gives
\[
 \tau_q(w)=\frac1\omega\log\frac{w_0}{w}
       +\int_w^\infty\left(\frac1{B(r)}-\frac1{\omega r}\right)\dd r+o(1).
\]
Since \(w_0e^{-\omega T_A}=h\), the time difference is
\(\omega^{-1}\log(h/H(w))+o(1)\). Bracketing \(H^{-1}(h)\) by two
nearby hit levels proves \eqref{an04cfsv1:positive} for \(w\), hence
for \(\arctan(qw)/q\). For the parameter compacts, \(h\) stays
bounded above and away from zero and the inverse levels stay a positive
distance above their respective roots. This follows by continuity of
the quadrature away from the root, its uniform logarithmic divergence
at the simple roots, and its uniform tail estimate. The same bracketing
is therefore uniform.

For the weight, \(\cos\theta\) stays bounded below on this portion of
the trajectory by the angular motion just proved. PROFILE coordinate
and \(K(t)=t+K(-t)\) give
\[
 (\log\widehat U_1)'=\tfrac12+d_q+
 \frac{q}{2\cos\theta}K(-\cos\theta/q),\qquad 0\le d_q\le q^2 .
\]
The integrated correction is \(o(1)\) since \(T_A=O(\ell)\).
At the endpoint \(\cos\theta=1+O(q^2)\); use
\(\widehat m=\widehat U_1^2/\cos^2\theta\).
Changing variables \(h=H(x)\) gives the density claimed.

For both endpoint strips, a bound independent of initial angle is
essential. In the positive-noise model \(m'\le(1+2q^2)m\), by bounding
the gated data second moment by its ungated matrix. Thus \(m(T_A)\le
Ae^{2q^2T_A}\); the zero-noise bound is \(A\). Each strip of initial
length \(\eta\) costs \(O(A\eta)\). First integrate on an interior
compact, then send \(\eta\downarrow0\). The endpoint labels themselves
have zero initial and terminal mass under this law, even though a
uniform per-label approximation fails there. The limiting gate boundary
has zero mass since both limiting maps are strictly increasing.
On Q1 the exact probe gate is \(\sin(q\zeta-\theta)>0\), equivalent
to \(x<\zeta\) for bounded offsets and small q. Integration gives the
fractions including equality at the support edge. Uniformity follows
also by taking convergent subsequences of the compact parameters and
offset: the label indicators converge away from at most one interior
boundary, with the same strip bound.
\end{proof}

In particular the requested positive-noise ratio \(x/h=1+o(1)\)
would require \(H^{-1}(h)=h\). In fact \(H^{-1}(h)>h\) for every
finite \(h>0\), since \(K(-\rho w)>0\).
The difference persists for a single fixed interior label. This is
the exact failure of the zero-noise substitution at the critical scale.

\paragraph{Other sectors.}
Q4 is excluded from both numerator and denominator. In the zero-noise
model its interior has \(\tan\theta=e^{-t/2}\tan\beta\), hence
\(x(T_A)\to0\) on compact interiors since \(1/(2\omega)>1\);
its mass also tends to \(A/8\) in total. No positive-noise Q4 profile
is used. Compact interior Q2 remains left-facing by CRA outer, with
gap of order \(q^{\lambda/(2\omega)}\) at \(T_A\). Its endpoint layers,
and all Q2 and Q3 labels, are excluded from the Q1 fractions.

\section{Saturation in the defined leading system}\label{an04cfsv1:leading}
Let a compact strong measure of mass \(m_0\in(0,1/2)\) start at \(t_0\)
in \([-R_0,R_0]^2\). Set the weak measure to zero, so \(N=V=K_w=0\).
In WIT's equations write \(\mathcal Y=\int y\,d\mu_s\),
\(\mathcal S(x)=\int F(\rho x,\rho\widetilde x)d\mu_s\).
With prime denoting normalized time the full system is
\begin{align}
 M'&=M(1-M),\qquad x'_1=-(1-M)x/2,\notag\\
 2x'_2&=\rho\phi(\rho x)-\rho\mathcal S(x)+\lambda y\Phi(\rho x),
 \qquad 2y'_2=\rho K(\rho x),\notag\\
 2y'_1&=-(1-M)y-\mathcal Y .
 \label{an04cfsv1:field}
\end{align}
Common radial reaction preserves normalized weights.
\begin{lemma}\label{an04cfsv1:sat}
The source-1-only input evolution is exactly \(x^{[1]}(t)
=x_0\sqrt{m_0/M(t)}\). The first half-mass time is \(t_0+D_0\),
\(D_0=\log[(1-m_0)/m_0]\). With
\(R_*=(1+R_0)e^{2D_0}-1\), the full input and output obey,
for \(0\le v=t-t_0\le D_0\),
\begin{align*}
 |x(t)-x_0\sqrt{m_0/M(t)}|
 &\le v\left[\tfrac34\rho\phi(0)+\tfrac34\lambda R_*\right],\\
 |y(t)-y_0\sqrt{m_0/M(t)}|
 &\le v\left[R_*/4+\tfrac\rho2(\phi(0)+\rho R_*)\right].
\end{align*}
\end{lemma}
\begin{proof}
Integrate \((1-M)=(\log M)'\) and the logistic equation.
The source-1-only formula for x does not solve the full x equation.
For an instantaneous radius \(R\), the Gaussian kernel satisfies
\(0\le F(t,r)\le K(r)\le\phi(0)+r_+\).
Thus \(|\mathcal S|\le M(\phi(0)+\rho R)\) and
\(|\mathcal Y|\le MR\). With \(M\le1/2\) these give
\(|x'_2|\le3\rho\phi(0)/4+3\lambda R/4\) and
\(|y'_2|\le\rho(\phi(0)+\rho R)/2\). The full speed in each coordinate
is at most \(2(1+R)\), so comparison of the support supremum yields
\(R\le R_*\). These bounds also give continuation to the stated finite
time by the compact-chart local Lipschitz characteristic equations.
Multiply the full equations by \(\sqrt M\) and integrate; the resulting
kernel \(\sqrt{M(u)/M(t)}\le1\) yields both bounds, including the
output donor term. In particular the isolated input multiplier at
half mass is \(\sqrt{2m_0}\).
\end{proof}

\begin{theorem}[Finite spreads and local rates]\label{an04cfsv1:spread}
On any finite chart-valid interval of the same \(N=0\) leading flow,
a pair with \(x_2\ge x_1,\ y_2\ge y_1\) at entry preserves that order.
Its output gap satisfies
\[
 y_2(t)-y_1(t)\ge\sqrt{m_0/M(t)}\,[y_2(t_0)-y_1(t_0)] .
\]
Thus an aligned nondegenerate entry profile of fixed positive mass
has a persistent scaled spread even as \(M\uparrow1\). On the exactly
saturated surface \(M=1\), its ordered output gap increases whenever
its input gap is positive. At the coherent resident the centered
symmetric mode grows at rate \(+\alpha_{\rm sh}\); the centered
antisymmetric mode contracts at rate \(-\alpha_{\rm sh}\).
\end{theorem}
\begin{proof}
For two actual receivers in one field write \(h=x_2-x_1,\ k=y_2-y_1\).
Their difference satisfies a linear equation with receiver Jacobian
averaged on the segment between them, with the current donors fixed.
Both off-diagonal entries are averages of \(\lambda\Phi(\rho x)/2>0\).
The diagonal entries are bounded on bounded charts. The integrating
factor for each diagonal, or the positive linear propagator, preserves
the nonnegative quadrant, including faces with one zero gap.
In particular
\[
 (\sqrt M\,k)'=\tfrac{\rho}{2}\sqrt M\,[K(\rho x_2)-K(\rho x_1)]\ge0 .
\]
The common output donor term cancels exactly for this pair.
Strict increase at \(M=1\) follows from strict monotonicity of K.
Full fields have at most linear growth in the support radius for
\(M\le1\), so fixed compact initial support continues on every fixed
finite interval. The bound also separates two weight quantiles with
an initial output gap, because the normalized label weights are
constant. It does not require a coherent donor field.

For the infinitesimal calculation write \(a=\rho K(\rho a)\).
P's \texttt{spectra} has \(F_t(t,t)=0,\ F_r(t,t)=\Phi(t)\);
centered perturbations have zero first-order donor variation. At
the resident the receiver matrix is
\[
 -\tfrac12(1-M)I+
 \begin{pmatrix}0&\alpha_{\rm sh}\\ \alpha_{\rm sh}&0\end{pmatrix}.
\]
This yields the asserted rates at \(M=1\), and rates \(-d\) and
\(-1/2-\alpha_{\rm sh}\) at \(M=0\). In contrast \(-\omega\)
is the strong outer zero-noise angular rate. The coherent mean block
at saturation is
\(\left(\begin{smallmatrix}-\alpha_{\rm sh}&\alpha_{\rm sh}\\
\alpha_{\rm sh}&-1/2\end{smallmatrix}\right)\), with negative trace
and determinant \(\alpha_{\rm sh}(1/2-\alpha_{\rm sh})>0\).
Mean attraction therefore coexists with shape instability.
\end{proof}
For the intended aligned spread the finite-amplitude conclusion is
persistence, with strict growth of the output gap at exact saturation.
A global exponential rate for arbitrary finite spreads does not follow
from the resident linearization. A special perturbation in the contracting
eigendirection has a different local behavior. The bounds here are for
the zero-weak-mass leading evolution; a growing weak measure changes
the field. If the separate weak-entry law \(N_0\asymp S_0^{1-\lambda}\)
and LOC's local hypotheses held, the nominal local gain through fixed
weak mass would scale as \(N_0^{-\alpha_{\rm sh}/\lambda}\asymp q^{-P}\)
for fixed c. This predicts loss of local smallness at weak half mass;
it is neither a lower bound on the actual spread there nor a valid
linear propagation of the full non-small profile. The limiting Q1
measures of Section~\ref{an04cfsv1:profiles} have unbounded scaled
support. Lemma~\ref{an04cfsv1:sat} applies to a prescribed compact
entry state, not automatically to an untruncated limit or to a truncated
subpopulation in the original self-consistent field.

\section{Trained-flow audit and S7 handoff}\label{an04cfsv1:original}
Fix \(0<\delta\le1/100\) before q. For small q, \(S_0<\delta\).
P's early envelope \(\bar S=S_0e^{(1+2q^2)t}\) and alignment
control apply through \(T_\delta=\log(\delta/S_0)/(1+2q^2)\);
\(q^2\log(\delta/S_0)\to0\). TAIL weightedmass requires alignment
error \(<\pi\), supplied there. TAIL earlyfraction gives
right-half ancestry mass at least
\(\delta e^{-4\delta}\exp[-2q^2\log(\delta/S_0)/(1+2q^2)]/4\).
ARR with \(R=1\) gives its distinct incoming Q2 sector mass
at least \(qS_0(S_0/\delta)^{1/16}/3000\). These are ancestry
budgets at the early clock, without a current scaled profile.
\begin{itemize}
\item First comparison failure in CRA's chain order: TAIL
\texttt{exponent} needs \((S_0/A)^\omega/q\to0\); the ratio
is exactly \(b>0\). The exact point-mass-data profile above replaces
that asymptotic.
\item First trained-profile transfer gap: an absolute early angle
allowance \(O(\delta)\) at fixed cap does not give \(o(q)\)
resolution on the critical input scale. At the nominal clock the
mass envelope is \(Ae^{2q^2T_A}\), of order one. No original labelwise
error or relative weighted-profile estimate of the necessary precision
is supplied.
\item PROFILE entry needs \(T\ge (2/\omega)\ell+C_1\);
the available clock is \(\ell/\omega+\log(A/c)\).
The endpoint \(\beta=\pi/2\) has the longer arrival coefficient.
\item OWPT's proposed substitution gives
\(\sigma=\kappa-2/\omega=-1/\omega\) at the limiting exponent;
exactly \(\sigma_q^{\rm pow}=-1/\omega-\log c/\ell\).
At its fixed overlap time the mass envelope has scale
\(c q^{-1/\omega}\) up to bounded factors. Neither its required
small mass nor its \(O(q^{\sigma-1})\) scaled coupling estimate is
obtained from that envelope. This is a failure of the bound, not a
lower bound for trained mass.
\item SMCP's fixed-offset passage requires the preceding overlap
and already found nonuniform outer-label distortion at ten.
MOMC's moving offset \(\ell/2\) additionally needs positive
\(\sigma-1/2,\ \sigma-3/2\) and small initial \(M_0\) before
half mass. Both powers have the wrong sign here. The deep time
\(10\ell-\log A_q+o(1)\) and its chart law cannot be used as a
critical-family initialized handoff.
\end{itemize}
For S7 the specified state is a pair of comparison Q1 measures and,
separately, a leading evolution estimate given compact ordered entry.
What is still needed is a trained entry clock, actual joint input/output
and radial profiles including the other sectors, tail field and rate
budgets, and transfer through saturation in the presence of their
backreaction. A positive inactive comparison fraction cannot certify a
trained bridge, its signed lag, or its producer-resolved effect.
\end{document}
